\chapter{广义函数}
\label{chap:distributions-fourier}

\section{广义函数的概念}
\label{sec:distribution-concept}

上一章中已经出现了连续谱的``本征态'' \(|x\rangle,|p\rangle\)。它们通常不能作为平方可积函数留在原 Hilbert 空间中；另一方面，Green 函数、点源和 Fourier 变换又不断用到 \(\delta\) 函数。为了把这些对象放在同一套严格语言中，引入广义函数（distribution）。

它的基本思想是：与其追问一个奇异对象在每一点``取什么值''，不如只考察它对一类足够光滑的检验函数产生什么作用。

先记

\[
\mathcal D(\mathbb R^n)\equiv C_c^\infty(\mathbb R^n),
\]

即 \(\mathbb R^n\) 上所有无穷可微并具有紧支集的复值函数组成的空间。这里下标 \(c\) 表示 compact support。

对普通函数 \(\phi\)，支集定义为

\[
\operatorname{supp}\phi
=\overline{\{\mathbf x\in\mathbb R^n:\phi(\mathbf x)\neq0\}}.
\]

因此 \(\phi\in C_c^\infty\) 的意思是：存在某个紧集 \(K\subset\mathbb R^n\)，使 \(\phi\) 在 \(K\) 外恒为零。

一个常用的球对称检验函数是

\[
\phi(\mathbf x)=
\begin{cases}
\exp\!\left[-\dfrac1{1-\|\mathbf x\|^2}\right],&\|\mathbf x\|<1,\\[4pt]
0,&\|\mathbf x\|\ge1.
\end{cases}
\]

它在球面 \(\|\mathbf x\|=1\) 上不但函数值为零，而且所有阶导数都能与外部的零函数光滑拼接。

为了简化高阶偏导，取多重指标

\[
\alpha=(\alpha_1,\ldots,\alpha_n),\qquad
|\alpha|=\sum_{j=1}^n\alpha_j,
\]

并记

\[
D^\alpha
=\frac{\partial^{|\alpha|}}
{\partial x_1^{\alpha_1}\cdots\partial x_n^{\alpha_n}}.
\]

于是一般的 \(m\) 阶线性微分算符可写为

\[
L=\sum_{|\alpha|\le m}a_\alpha(\mathbf x)D^\alpha.
\]

这里 \(a_\alpha\) 是已知系数函数。

检验函数空间上的收敛不能只看函数值。称 \(\phi_j\to\phi\) 于 \(\mathcal D(\mathbb R^n)\)，若存在一个共同的紧集 \(K\)，使

\[
\operatorname{supp}\phi_j\subset K,
\qquad
\operatorname{supp}\phi\subset K,
\]

并且对每一个多重指标 \(\alpha\) 都有

\[
\sup_{\mathbf x\in K}
|D^\alpha\phi_j(\mathbf x)-D^\alpha\phi(\mathbf x)|\to0.
\]

也就是说，在一个共同的紧集上，函数以及任意阶导数都一致收敛。

现在可以给出定义。一个线性泛函

\[
T:\mathcal D(\mathbb R^n)\to\mathbb C
\]

称为广义函数或分布，若它关于上述拓扑连续。通常把它对检验函数的作用记为

\[
\langle T,\phi\rangle.
\]

线性意味着

\[
\langle T,a\phi+b\psi\rangle
=a\langle T,\phi\rangle+b\langle T,\psi\rangle,
\]

连续性意味着只要 \(\phi_j\to\phi\) 于 \(\mathcal D\)，就有

\[
\langle T,\phi_j\rangle\to\langle T,\phi\rangle.
\]

全体分布组成的空间记为

\[
\mathcal D'(\mathbb R^n).
\]

最重要的第一个例子是 Dirac \(\delta\) 分布。对固定 \(\boldsymbol\xi\in\mathbb R^n\)，定义

\[
\boxed{
\langle\delta(\mathbf x-\boldsymbol\xi),\phi\rangle
=\phi(\boldsymbol\xi)
}.
\]

这一定义并不是说 \(\delta\) 是某个``在一点无穷大、其余地方为零''的普通函数；它只规定了 \(\delta\) 对任意检验函数的作用。

再看普通函数怎样嵌入分布空间。若 \(f\in L^1_{\mathrm{loc}}(\mathbb R^n)\)，即对任意紧集 \(K\) 都有

\[
\int_K|f(\mathbf x)|\,\dd^n x<\infty,
\]

则可定义

\[
\boxed{
\langle T_f,\phi\rangle
=\int_{\mathbb R^n}f(\mathbf x)\phi(\mathbf x)\,\dd^n x
}.
\]

因为 \(\phi\) 有紧支集，积分只发生在有限区域内。这样的分布称为正则分布。若两个局部可积函数几乎处处相等，它们给出同一个正则分布，所以这里真正对应的是 \(L^1_{\mathrm{loc}}\) 中按几乎处处相等识别的等价类。

这里的``连续''也可以直接核对。设 \(f\in L^1_{\mathrm{loc}}\)，并且 \(\phi_n\to0\) 于 \(\mathcal D\)。按 \(\mathcal D\) 的收敛定义，存在同一个紧集 \(K\) 包含所有 \(\operatorname{supp}\phi_n\)，并且

\[
\|\phi_n\|_\infty\to0.
\]

于是

\[
|\langle T_f,\phi_n\rangle|
\le\int_K|f(x)|\,|\phi_n(x)|\,\dd x
\le\|\phi_n\|_\infty\int_K|f(x)|\,\dd x\to0.
\]

所以局部可积函数并不是``形式上''嵌入分布空间，而是确实给出了连续线性泛函。

不是每一个``点态写出的函数''都自动给出正则分布；要由上面的 Lebesgue 积分定义正则分布，至少需要局部可积。这里不展开完整测度论，只使用以下事实：Lebesgue 可测集构成一个 \(\sigma\)-代数，而 \(L^1_{\mathrm{loc}}\) 中的函数均按 Lebesgue 积分理解。

不是由局部可积函数直接产生的分布称为奇异分布，\(\delta\) 及其导数就是典型例子。

分布一般没有逐点值，但仍然可以定义局部为零与支集。若开集 \(U\) 中所有检验函数都满足

\[
\operatorname{supp}\phi\subset U
\quad\Longrightarrow\quad
\langle T,\phi\rangle=0,
\]

就称 \(T\) 在 \(U\) 上为零。\(T\) 的支集定义为它``不在任何邻域上恒为零''的点所组成的闭集。特别地，

\[
\operatorname{supp}\delta(\mathbf x-\boldsymbol\xi)=\{\boldsymbol\xi\}.
\]


接下来考察局部正则化、奇异支集与 mollifier。

分布的支集只回答``对象在哪里不为零''，却不能区分一个点附近究竟是光滑函数、跳跃还是 delta 型奇性。为描述正则性，先引入标准 mollifier。取
\[
\rho\in C_c^\infty(\RR^n),
\qquad
\rho\ge0,
\qquad
\int_{\RR^n}\rho(x)\,\dd^n x=1,
\]
并令
\[
\rho_\varepsilon(x)=\varepsilon^{-n}\rho(x/\varepsilon).
\]
对分布 $T$，卷积 $T*\rho_\varepsilon$ 是光滑函数；若 $T=T_f$ 来自 $f\in L^p_{\mathrm{loc}}$，则在适当的局部意义下
\[
T*\rho_\varepsilon\longrightarrow T,
\qquad
f*\rho_\varepsilon\longrightarrow f.
\]
这给出一个重要原则：分布并不是``无法计算的奇异对象''，而是可以被一族光滑对象稳定逼近；真正的奇性体现在当 $\varepsilon\to0$ 时某些导数不再一致有界。

一个开集 $U$ 上若存在 $f\in C^\infty(U)$ 使
\[
\langle T,\phi\rangle
=\int_U f(x)\phi(x)\,\dd^n x,
\qquad
\forall\phi\in C_c^\infty(U),
\]
就称 $T$ 在 $U$ 上光滑。于是定义奇异支集
\[
\boxed{
\operatorname{singsupp}T
:=
\RR^n\setminus
\bigcup\{U:T|_U\in C^\infty(U)\}
}.
\]
总有
\[
\operatorname{singsupp}T\subseteq\operatorname{supp}T,
\]
但二者通常相差很大。例如 Heaviside 分布 $H(x)$ 的支集是 $[0,\infty)$，而
\[
\operatorname{singsupp}H=\{0\}.
\]
同理，$\delta$ 与其任意阶导数的奇异支集都只是原点。

mollifier 还能把分布微分与经典微分连接起来。因为卷积与常系数微分交换，
\[
D^\alpha(T*\rho_\varepsilon)
=(D^\alpha T)*\rho_\varepsilon
=T*(D^\alpha\rho_\varepsilon).
\]
所以可以先在光滑层面完成积分分部、PDE 估计和极限，再回到分布。这是弱解理论中``先平滑、后估计、再取极限''的标准逻辑。

接下来考察波前集：奇性的位置与方向。

奇异支集仍只记录位置，无法区分一个奇点沿哪些共变方向缺乏正则性。Fourier 变换提供了方向信息：若紧支分布 $u$ 实际上是 $C^\infty$ 函数，则其 Fourier 变换在任意方向都比任何幂更快衰减。反过来，某个锥形方向上缺乏快速衰减就意味着存在该方向的高频奇性。

取 $x_0\in\RR^n$、$\xi_0\in\RR^n\setminus\{0\}$。若存在 $\chi\in C_c^\infty$，满足 $\chi(x_0)\neq0$，并存在包含 $\xi_0$ 的开锥 $\Gamma$，使对任意 $N\ge0$ 都存在 $C_N$ 满足
\[
\boxed{
|\widehat{\chi T}(\xi)|
\le C_N(1+|\xi|)^{-N},
\qquad \xi\in\Gamma
},
\]
则称 $T$ 在 $(x_0,\xi_0)$ 处 microlocally smooth。所有不满足这一条件的点组成波前集
\[
\operatorname{WF}(T)
\subset T^*\RR^n\setminus0.
\]
这里去掉零截面，是因为 $\xi=0$ 不代表高频方向。波前集在纤维变量中是锥形的：
\[
(x,\xi)\in\operatorname{WF}(T)
\Longrightarrow
(x,\lambda\xi)\in\operatorname{WF}(T),
\qquad \lambda>0.
\]
并且把所有方向忘掉后正好恢复奇异支集：
\[
\boxed{
\operatorname{singsupp}T
=
\pi\bigl(\operatorname{WF}(T)\bigr)
},
\]
其中 $\pi:T^*\RR^n\to\RR^n$ 是余切丛投影。

最简单的例子是 $\delta_0$。因为任取 $\chi(0)\neq0$，有
\[
\widehat{\chi\delta_0}(\xi)=\chi(0),
\]
在任何非零方向都不衰减，所以
\[
\operatorname{WF}(\delta_0)
=\{(0,\xi):\xi\neq0\}.
\]
相比之下，对一维 Heaviside 分布，奇性仍位于 $x=0$，且两个非零共变方向都出现，因此
\[
\operatorname{WF}(H)
=\{(0,\xi):\xi\neq0\}.
\]
高维曲面 delta 分布则更有信息：若 $\Sigma=\{x:F(x)=0\}$ 且 $\dd F\neq0$，则 $\delta(F)$ 的波前方向沿曲面的 conormal bundle，即只出现在法向余切方向，而不是任意切向方向。

波前集最重要的用途之一是判断两个分布何时可以相乘。一般分布乘法并不存在；例如形式上写 $\delta^2$ 没有 canonical 意义。Hörmander 的乘积判据给出一个精确充分条件：若不存在 $x$ 和非零 $\xi$ 使
\[
(x,\xi)\in\operatorname{WF}(u),
\qquad
(x,-\xi)\in\operatorname{WF}(v),
\]
则乘积 $uv$ 可以 canonical 地定义为分布。直观上，两个因子的高频奇性不能在同一点以相反余切方向正面对撞；否则卷积型 Fourier 表达式会失去可控性。

\begin{note}[为什么本章只建立最低 microlocal 框架]
波前集真正强大的地方在于研究偏微分算子的奇性传播：椭圆算子不会制造新的波前方向，双曲算子则沿特征流传播奇性。完整理论需要 pseudodifferential operator、symbol calculus 与 Hamilton flow。本书在这里先固定定义和乘积判据；后续若需要传播定理，再从这一母结构继续展开，而不把 microlocal analysis 退化成一张术语表。
\end{note}

\section{广义函数的基本运算}
\label{sec:distribution-operations}

分布的线性运算按对检验函数的作用定义：

\[
\langle aT+bS,\phi\rangle
=a\langle T,\phi\rangle+b\langle S,\phi\rangle.
\]

若 \(a(\mathbf x)\in C^\infty(\mathbb R^n)\)，则 \(a\phi\) 仍属于 \(C_c^\infty\)，因此可以定义光滑函数与分布的乘积

\[
\boxed{
\langle aT,\phi\rangle
=\langle T,a\phi\rangle
}.
\]

这与正则分布的普通乘法完全相容。

分布最重要的运算之一是微分。对一维分布定义

\begin{equation}
\boxed{
\langle T',\phi\rangle=-\langle T,\phi'\rangle
}.
\label{eq:distribution-derivative}
\end{equation}

多维情形为

\[
\boxed{
\langle D^\alpha T,\phi\rangle
=(-1)^{|\alpha|}\langle T,D^\alpha\phi\rangle
}.
\]

这个定义来自普通函数的分部积分。由于检验函数具有紧支集，边界项自动消失，因此即使 \(T\) 本身没有经典导数，分布导数仍然存在。

例如

\[
\langle\delta^{(\alpha)}(\mathbf x-\boldsymbol\xi),\phi\rangle
=(-1)^{|\alpha|}D^\alpha\phi(\boldsymbol\xi).
\]

再看一个有跳跃的分段光滑函数：

\[
f(x)=
\begin{cases}
f_1(x),&x<x_0,\\
f_2(x),&x>x_0,
\end{cases}
\]

其中两侧都连续可微。由分布导数定义，

\[
\begin{aligned}
\langle f',\phi\rangle
&=-\int_{-\infty}^{x_0}f_1\phi'\,\dd x
  -\int_{x_0}^{\infty}f_2\phi'\,\dd x\\
&=\int_{-\infty}^{x_0}f_1'\phi\,\dd x
  +\int_{x_0}^{\infty}f_2'\phi\,\dd x
  +[f_2(x_0)-f_1(x_0)]\phi(x_0).
\end{aligned}
\]

所以

\[
\boxed{
f'=f_c'+[f_2(x_0)-f_1(x_0)]\delta(x-x_0)
}.
\]

特别地，Heaviside 函数

\[
H(x)=
\begin{cases}0,&x<0,\\1,&x>0\end{cases}
\]

满足

\[
\boxed{H'=\delta}.
\]

另一个常见例子是

\[
\boxed{
\frac{\dd}{\dd x}\ln|x|
=\operatorname{pv}\frac1x
},
\]

其中右端是 Cauchy 主值分布，而不是普通意义下在 \(x=0\) 可积的函数。

变量代换也要通过检验函数来定义。设

\[
T:\mathbb R^n\to\mathbb R^n,
\qquad \mathbf y=T(\mathbf x)
\]

是 \(C^\infty\) 微分同胚。对分布 \(F\) 定义其拉回 \(F\circ T\) 为

\[
\boxed{
\langle F\circ T,\phi\rangle
=
\left\langle
F(\mathbf y),
\left|\det DT^{-1}(\mathbf y)\right|
\phi\!\left(T^{-1}(\mathbf y)\right)
\right\rangle
}.
\]

若 \(F\) 恰好由局部可积函数产生，这个式子就是普通多元积分换元公式。

一维情况下，若 \(g(x)\) 的零点 \(x_j\) 都是单根，

\[
g(x_j)=0,
\qquad g'(x_j)\neq0,
\]

则由局部换元得到

\[
\boxed{
\delta(g(x))
=\sum_j\frac{\delta(x-x_j)}{|g'(x_j)|}
}.
\]

验证时对任意检验函数 \(\phi\) 有

\[
\langle\delta(g(x)),\phi\rangle
=\sum_j\frac{\phi(x_j)}{|g'(x_j)|}.
\]

这个系数可以从每个零点附近的普通换元看出来。由于 \(x_j\) 是单根，逆函数定理保证在一个足够小的邻域 \(U_j\) 内，\(y=g(x)\) 是微分同胚。只看支集落在 \(U_j\) 内的检验函数，有

\[
\begin{aligned}
\langle\delta(g(x)),\phi(x)\rangle
&=\left\langle\delta(y),
\phi(x(y))\left|\frac{\dd x}{\dd y}\right|\right\rangle\\
&=\phi(x_j)\left|\frac1{g'(x_j)}\right|.
\end{aligned}
\]

再用单位分解把一般检验函数拆成各零点邻域中的部分并求和，就得到上式。绝对值不能省略，因为 delta 分布在变量反向时仍应保持正的测度权重。

因此尺度变换给出

\[
\delta(\alpha x)=\frac1{|\alpha|}\delta(x),
\]

而在 \(n\) 维中

\[
\boxed{
\delta^{(n)}(\alpha\mathbf x)
=\frac1{|\alpha|^n}\delta^{(n)}(\mathbf x)
}.
\]

若用一组新坐标 \(\mathbf u\) 参数化 \(\mathbf x=\mathbf x(\mathbf u)\)，并且在所讨论区域内变换可逆，记

\[
J(\mathbf u)=\det\frac{\partial\mathbf x}{\partial\mathbf u},
\]

则体积元满足

\[
\dd^n x=|J(\mathbf u)|\,\dd^n u.
\]

相应地，若 \(\mathbf x(\mathbf u')=\mathbf x'\)，则

\[
\boxed{
\delta^{(n)}(\mathbf x(\mathbf u)-\mathbf x')
=\frac{\delta^{(n)}(\mathbf u-\mathbf u')}{|J(\mathbf u')|}
}.
\]

这就是曲线坐标中 \(\delta\) 函数必须带 Jacobian 的原因。

分布还可以取原函数。若存在分布 \(G\) 使

另一个反复出现的对象是 Cauchy 主值分布。定义

\begin{equation}
\boxed{
\left\langle\operatorname{pv}\frac1x,\phi\right\rangle
=\lim_{\varepsilon\downarrow0}
\int_{|x|>\varepsilon}\frac{\phi(x)}x\,\dd x
}.
\label{eq:pv-definition}
\end{equation}

把正、负半轴配对后，等价地写成

\[
\left\langle\operatorname{pv}\frac1x,\phi\right\rangle
=\int_0^\infty\frac{\phi(x)-\phi(-x)}x\,\dd x,
\]

右端在 \(x=0\) 附近是良性的，因为

\[
\phi(x)-\phi(-x)=2x\phi'(0)+O(x^3).
\]

这也说明普通函数 \(1/x\) 虽然在零点不局部可积，但通过对称消去仍可定义一个分布。

对 \(\ln|x|\) 使用分布导数定义，

\[
\left\langle(\ln|x|)',\phi\right\rangle
=-\int_{-\infty}^{\infty}\ln|x|\,\phi'(x)\,\dd x.
\]

在 \((-\infty,-\varepsilon)\) 与 \((\varepsilon,\infty)\) 上分别分部积分后，边界的 \(\ln\varepsilon\) 项正好相消，令 \(\varepsilon\downarrow0\) 得

\[
\boxed{(\ln|x|)'=\operatorname{pv}\frac1x}.
\]

这就是下面边界值公式中 \(1/x\) 必须理解为主值分布的原因。

\[
G'=F,
\]

就称 \(G\) 为 \(F\) 的一个原函数。它和普通不定积一样，只确定到一个常数分布。

边界值分布经常写成

\[
\ln(x\pm \ii 0)
=\ln|x|\pm \ii\pi H(-x).
\]

对它求分布导数得到

\[
\boxed{
\frac{\dd}{\dd x}\ln(x+\ii 0)
=\operatorname{pv}\frac1x-\ii\pi\delta(x)
},
\]

\[
\boxed{
\frac{\dd}{\dd x}\ln(x-\ii 0)
=\operatorname{pv}\frac1x+\ii\pi\delta(x)
}.
\]

最后，分布序列的收敛也按对所有检验函数的作用来定义。若对任意 \(\phi\in\mathcal D(\mathbb R^n)\) 都有

\[
\langle T_j,\phi\rangle\to\langle T,\phi\rangle,
\]

就记

\[
T_j\to T\qquad\text{于 }\mathcal D'(\mathbb R^n).
\]

这种``弱作用意义''的收敛允许许多普通函数序列以 \(\delta\) 或其导数为极限，也为下一节把 Fourier 变换扩充到缓增分布提供了自然语言。


接下来考察局部阶数与分布微分的连续性。

分布连续性的定量形式是后续估计中真正使用的版本。对每个紧集 \(K\subset\RR^n\)，若 \(T\in\mathcal D'(\RR^n)\)，则存在常数 \(C_K>0\) 与非负整数 \(m_K\)，使所有满足 \(\operatorname{supp}\phi\subset K\) 的检验函数都有
\begin{equation}
|\langle T,\phi\rangle|
\le
C_K
\max_{|\alpha|\le m_K}
\sup_{x\in K}|D^\alpha\phi(x)|.
\label{eq:distribution-local-order-bound}
\end{equation}
这说明一个 distribution 在每个固定紧集上只需要有限阶导数半范数就能控制；\(m_K\) 称为该紧集上的一个阶数上界。若可以对所有紧集取同一个 \(m\)，才称 \(T\) 具有全局有限阶。一般 distribution 不必有统一的全局阶数。

分布导数不会破坏这一结构。若\eqnrefs{eq:distribution-local-order-bound} 对某个 \(m_K\) 成立，则
\[
|\langle D^\beta T,\phi\rangle|
=
|\langle T,(-1)^{|\beta|}D^\beta\phi\rangle|
\le
C_K
\max_{|\alpha|\le m_K+|\beta|}
\sup_K|D^\alpha\phi|.
\]
因此 distribution derivative 始终仍是 distribution，而其局部阶数至多增加 \(|\beta|\)。

\begin{derivation}[分布导数与经典导数严格相容]
设 \(f\in C^1(\RR^n)\cap L^1_{\mathrm{loc}}(\RR^n)\)，令 \(T_f\) 为其正则分布，即
\[
\langle T_f,\phi\rangle=\int_{\RR^n}f(x)\phi(x)\,\dd^nx.
\]
对第 \(j\) 个坐标方向，分布导数的定义是对任意检验函数 \(\phi\in C_c^\infty(\RR^n)\) 令
\[
\langle \partial_jT_f,\phi\rangle
:=-\langle T_f,\partial_j\phi\rangle.
\]
把正则分布代入右端，
\[
\langle \partial_jT_f,\phi\rangle
=-\int_{\RR^n}f(x)\partial_j\phi(x)\,\dd^nx.
\]
因为 \(\phi\) 有紧支集 \(K=\operatorname{supp}\phi\)，取包含 \(K\) 的有界开集 \(\Omega\)，积分区域可换成 \(\Omega\)。由 Fubini 定理先对 \(x^j\) 积分：记 \(\dd\hat x^j=\dd x^1\cdots\widehat{\dd x^j}\cdots\dd x^n\)，则
\[
\int_\Omega f\,\partial_j\phi\,\dd^nx
=\int_{\RR^{n-1}}
\left(
\int_{-\infty}^{+\infty} f(\hat x^j,x^j)\,\partial_j\phi(\hat x^j,x^j)\,\dd x^j
\right)
\dd\hat x^j.
\]
对内层一维积分作经典分部积分（\(\partial_j\phi\) 关于 \(x^j\) 紧支集）：
\[
\int_{-\infty}^{+\infty} f\,\partial_j\phi\,\dd x^j
=\bigl[f\phi\bigr]_{-\infty}^{+\infty}
-\int_{-\infty}^{+\infty}(\partial_jf)\phi\,\dd x^j
=-\int_{-\infty}^{+\infty}(\partial_jf)\phi\,\dd x^j.
\]
回代即得 Green 恒等式
\[
\int_\Omega f\,\partial_j\phi\,\dd^nx
=-\int_\Omega(\partial_jf)\phi\,\dd^nx
+\int_{\partial\Omega}f\phi\,n_j\,\dd S.
\]
因 \(\phi\) 在 \(\partial\Omega\) 上为零，边界项消失，故
\[
-\int_{\RR^n} f\,\partial_j\phi\,\dd^nx
=\int_{\RR^n}(\partial_jf)\phi\,\dd^nx
=\langle T_{\partial_jf},\phi\rangle.
\]
上式对一切 \(\phi\in C_c^\infty(\RR^n)\) 成立，由分布相等判据得到
\[
\partial_jT_f=T_{\partial_jf}.
\]
因此 distribution derivative 不是重新定义了一个与经典微分无关的运算，而是经典导数的唯一连续延拓。

\emph{举例：Heaviside 函数产生 delta 分布。}取 \(n=1\)、\(H(x)=\mathbf 1_{x>0}\)。对任意 \(\phi\in C_c^\infty(\RR)\)，
\[
\langle H',\phi\rangle
=-\int_0^\infty\phi'(x)\,\dd x
=-\bigl[\phi(x)\bigr]_0^\infty
=\phi(0)
=\langle\delta_0,\phi\rangle.
\]
故 \(H'=\delta_0\)。这正是弱导数记录跳跃奇点的机制。

\emph{推广与物理建模。}对分段 \(C^1\) 函数 \(f\) 在超曲面 \(\Sigma\) 处有跳跃 \([f]_\Sigma\)，分布导数满足
\[
\partial_jT_f=T_{\partial_jf}+[f]_\Sigma\,n_j\,\delta_\Sigma,
\]
其中 \(n_j\) 为 \(\Sigma\) 的法向分量。量子力学中一维 delta 势阱 \(V(x)=-\lambda\delta(x)\) 的 Schrödinger 方程正是在分布意义下求解：波函数 \(\psi\) 连续但 \(\psi'\) 在 \(x=0\) 处有跳跃
\[
\psi'(0^+)-\psi'(0^-)=-\frac{2m\lambda}{\hbar^2}\psi(0),
\]
此即分布导数把 delta 奇性翻译为导数跳跃条件。高阶情形 \(\partial^\alpha T_f=T_{\partial^\alpha f}\) 由归纳法直接得到。
\end{derivation}

\section{缓增广义函数和 Fourier 变换}
\label{sec:tempered-fourier}

普通检验函数 \(C_c^\infty(\mathbb R^n)\) 适合定义一般分布，但 Fourier 变换会把紧支集函数变成通常不再紧支集的函数。为了让 Fourier 变换在检验函数空间内部闭合，引入 Schwartz 空间

\[
\mathcal S(\mathbb R^n).
\]

函数 \(\phi\in C^\infty(\mathbb R^n)\) 属于 \(\mathcal S(\mathbb R^n)\)，若对任意多重指标 \(\alpha,\beta\) 都有

\[
p_{\alpha\beta}(\phi)
=\sup_{x\in\mathbb R^n}|x^\alpha D^\beta\phi(x)|<\infty.
\]

这表示 \(\phi\) 及其所有导数都比任意负幂更快衰减。由所有半范数 \(p_{\alpha\beta}\) 定义拓扑后，\(\mathcal S\) 是完备局部凸空间，而且

\[
C_c^\infty(\mathbb R^n)\subset\mathcal S(\mathbb R^n)\subset L^2(\mathbb R^n)
\]

并且两次嵌入都是稠密的。

接下来考察Fourier 变换约定。

本讲义统一采用对称归一化

\begin{equation}
\widehat\phi(k)=\mathcal F[\phi](k)
=\frac1{(2\pi)^{n/2}}
\int_{\mathbb R^n}\phi(x)\ee^{-\ii k\cdot x}\,\dd^n x,
\label{eq:fourier-convention}
\end{equation}

\[
\phi(x)=\mathcal F^{-1}[\widehat\phi](x)
=\frac1{(2\pi)^{n/2}}
\int_{\mathbb R^n}\widehat\phi(k)\ee^{\ii k\cdot x}\,\dd^n k.
\]

Fourier 变换把 \(\mathcal S(\mathbb R^n)\) 连续双射到自身。最常用的两条计算规则来自分部积分和对参数求导：

\[
\mathcal F[D^\alpha f](k)=(ik)^\alpha\widehat f(k),
\]

\[
\mathcal F[x^\alpha f](k)=\ii^{|\alpha|}D_k^\alpha\widehat f(k).
\]

以一维一阶导数为例，

\[
\begin{aligned}
\mathcal F[f'](k)
&=\frac1{\sqrt{2\pi}}\int_{-\infty}^{\infty}f'(x)\ee^{-\ii kx}\,\dd x\\
&=\frac1{\sqrt{2\pi}}\Big[f(x)\ee^{-\ii kx}\Big]_{-\infty}^{\infty}
+ik\widehat f(k)\\
&=\ii k\widehat f(k),
\end{aligned}
\]

其中边界项为零正是 Schwartz 衰减条件的作用。

另外

\[
\mathcal F^2f(x)=f(-x),\qquad \mathcal F^4=I.
\]

这应称为 Fourier 变换的二次作用性质，而不是``共轭性质''。


Fourier 变换首先在 Schwartz 空间上封闭；这一点来自多项式权与微分在 Fourier 变换下互换。

上面使用了
\[
\mathcal F:\mathcal S(\mathbb R^n)\to\mathcal S(\mathbb R^n)
\]
这一事实。它之所以重要，是因为缓增分布的 Fourier 变换正是靠这个闭性定义的。这里把证明中的核心估计写出来。

设 $\phi\in\mathcal S(\mathbb R^n)$。由于 $\phi$ 比任意幂次都快地衰减，$x^\beta\phi(x)$ 对任意多重指标 $\beta$ 都属于 $L^1(\mathbb R^n)$。因此可以在积分号下对 $k$ 求导：
\[
D_k^\beta\widehat\phi(k)
=\frac1{(2\pi)^{n/2}}
\int_{\mathbb R^n}(-\ii x)^\beta
\phi(x)\ee^{-\ii k\cdot x}\,\dd^n x,
\]
也就是
\begin{equation}
D_k^\beta\widehat\phi
=(-\ii)^{|\beta|}\widehat{x^\beta\phi}.
\label{eq:schwartz-fourier-derivative}
\end{equation}
所以 $\widehat\phi$ 无穷可微。

还需要证明它的每个导数仍快速衰减。利用
\[
\mathcal F[D_x^\alpha h](k)
=(\ii k)^\alpha\widehat h(k),
\]
把 \cref{eq:schwartz-fourier-derivative} 乘以 $k^\alpha$，得到
\begin{equation}
k^\alpha D_k^\beta\widehat\phi(k)
=C_{\alpha\beta}
\mathcal F\!\left[D_x^\alpha(x^\beta\phi)\right](k),
\label{eq:schwartz-seminorm-transform}
\end{equation}
其中 $|C_{\alpha\beta}|=1$，具体相位对估计没有影响。对任意 $h\in L^1$，Fourier 定义直接给出
\[
\|\widehat h\|_\infty
\le\frac1{(2\pi)^{n/2}}\|h\|_1.
\]
因此
\[
\sup_k
\left|k^\alpha D_k^\beta\widehat\phi(k)\right|
\le
\frac1{(2\pi)^{n/2}}
\left\|D_x^\alpha(x^\beta\phi)\right\|_1.
\]
Leibniz 公式把右端展开成有限个
\[
x^\gamma D^\delta\phi(x)
\]
的 $L^1$ 范数。为了把 $L^1$ 范数继续控制到 Schwartz 半范数，只需取整数 $N>n$，写
\[
|x^\gamma D^\delta\phi(x)|
\le
(1+|x|)^{-N}
\sup_y(1+|y|)^N|y^\gamma D^\delta\phi(y)|.
\]
由于
\[
\int_{\mathbb R^n}(1+|x|)^{-N}\,\dd^n x<\infty,
\]
便得到
\[
\sup_k
\left|k^\alpha D_k^\beta\widehat\phi(k)\right|
\le C_{\alpha\beta N}
\sum_{(\gamma,\delta)\in F_{\alpha\beta N}}
 p_{\gamma\delta}(\phi),
\]
其中 $F_{\alpha\beta N}$ 只是一个有限指标集。于是对所有 $\alpha,\beta$，
\[
p_{\alpha\beta}(\widehat\phi)<\infty,
\]
即
\begin{equation}
\boxed{\widehat\phi\in\mathcal S(\mathbb R^n).}
\label{eq:fourier-schwartz-invariance}
\end{equation}
而且上面的有限半范数估计同时证明了 $\mathcal F$ 在 Schwartz 拓扑下连续。

对逆 Fourier 变换作完全相同的估计可知 $\mathcal F^{-1}$ 也连续地把 $\mathcal S$ 映到自身。结合后面用 Gaussian 核证明的 Fourier 反演，就得到
\[
\mathcal F:\mathcal S\longrightarrow\mathcal S
\]
是拓扑线性同构。于是对 $T\in\mathcal S'$ 定义
\[
\langle\mathcal FT,\phi\rangle
=\langle T,\mathcal F\phi\rangle
\]
时，右端始终有意义；这正是缓增分布 Fourier 理论成立的关键。


接下来考察缓增广义函数。

Schwartz 空间的连续线性泛函称为缓增广义函数，记为

\[
\mathcal S'(\mathbb R^n).
\]

普通函数若至多按多项式增长，并且局部可积，就能通过

\[
\langle T_f,\phi\rangle
=\int_{\mathbb R^n}f(x)\phi(x)\,\dd^n x
\]

定义一个缓增广义函数。但\textbf{并非每个缓增广义函数都能写成普通函数积分的形式}；\(\delta\) 及其导数就是最基本的反例。

缓增分布上的 Fourier 变换通过对 Schwartz 空间的对偶作用定义。

对 \(T\in\mathcal S'(\mathbb R^n)\)，定义

\[
\boxed{\langle\mathcal F T,\phi\rangle
=\langle T,\mathcal F\phi\rangle},
\qquad \phi\in\mathcal S(\mathbb R^n).
\]

由于 \(\mathcal F:\mathcal S\to\mathcal S\) 连续，这一定义确实给出新的缓增广义函数，并与普通函数 Fourier 变换相容。

在本归一化下，

\[
\mathcal F[1]=(2\pi)^{n/2}\delta(k),
\qquad
\mathcal F[\delta(x)]=\frac1{(2\pi)^{n/2}}.
\]

平移、调制与缩放为

\[
\mathcal F[f(x-x_0)](k)=\ee^{-\ii k\cdot x_0}\widehat f(k),
\]

\[
\mathcal F[\ee^{\ii k_0\cdot x}f(x)](k)=\widehat f(k-k_0),
\]

\[
\mathcal F[f(Ax)](k)=\frac1{|\det A|}\widehat f(A^{-T}k),
\]

其中 \(A\) 为可逆实矩阵。一维 \(A=\alpha\) 时即

\[
\mathcal F[f(\alpha x)](k)=\frac1{|\alpha|}\widehat f(k/\alpha).
\]


\begin{derivation}[tempered distribution 上 Fourier 微分律]
设 \(T\in\mathcal S'(\RR^n)\)。按本章约定
\[
\langle\mathcal FT,\phi\rangle
=\langle T,\mathcal F\phi\rangle.
\]
对任意 \(\phi\in\mathcal S\)，先计算
\[
D_x^\alpha\mathcal F\phi(x)
=
\mathcal F\bigl[(-\ii k)^\alpha\phi(k)\bigr](x).
\]
于是
\[
\begin{aligned}
\langle\mathcal F[D^\alpha T],\phi\rangle
&=\langle D^\alpha T,\mathcal F\phi\rangle\\
&=(-1)^{|\alpha|}
\langle T,D_x^\alpha\mathcal F\phi\rangle\\
&=(-1)^{|\alpha|}
\langle T,\mathcal F[(-\ii k)^\alpha\phi]\rangle\\
&=\langle\mathcal FT,(\ii k)^\alpha\phi\rangle.
\end{aligned}
\]
因此作为 tempered distributions，
\begin{equation}
\mathcal F[D^\alpha T]
=(\ii k)^\alpha\mathcal FT.
\label{eq:tempered-fourier-derivative}
\end{equation}
同理，由
\[
\mathcal F[k^\alpha\phi](x)
=\ii^{|\alpha|}D_x^\alpha\mathcal F\phi(x)
\]
可得
\begin{equation}
\mathcal F[x^\alpha T]
=\ii^{|\alpha|}D_k^\alpha\mathcal FT.
\label{eq:tempered-fourier-multiplication}
\end{equation}
这两式是用 Fourier 方法构造常系数 PDE 基本解的代数核心。
\end{derivation}

接下来考察Parseval / Plancherel 关系。

对 \(f,g\in\mathcal S\)，有

\[
\int_{\mathbb R^n}f^*(x)g(x)\,\dd^n x
=\int_{\mathbb R^n}\widehat f^{\,*}(k)\widehat g(k)\,\dd^n k.
\]

取 \(g=f\) 得

\[
\|f\|_{L^2}=\|\widehat f\|_{L^2}.
\]

由 \(\mathcal S\) 在 \(L^2\) 中稠密，可把 Fourier 变换唯一延拓为 \(L^2\) 上的幺正算符。这正是\chapref{chap:hilbert-space} Hilbert 空间结构在 Fourier 分析中的具体体现。

在 Schwartz 函数上，Parseval 关系可以直接由 Fourier 反演算出。由

\[
g(x)=\frac1{(2\pi)^{n/2}}
\int\widehat g(k)\ee^{\ii k\cdot x}\,\dd^nk,
\]

有

\[
\begin{aligned}
\langle f|g\rangle
&=\int f^*(x)g(x)\,\dd^nx\\
&=\frac1{(2\pi)^{n/2}}
\int\!\int f^*(x)\widehat g(k)\ee^{\ii k\cdot x}\,\dd^nk\,\dd^nx.
\end{aligned}
\]

Schwartz 衰减保证可以交换积分次序，而内层积分正是

\[
\int f^*(x)\ee^{\ii k\cdot x}\,\dd^nx
=(2\pi)^{n/2}\widehat f^{\,*}(k).
\]

因此

\[
\langle f|g\rangle
=\int\widehat f^{\,*}(k)\widehat g(k)\,\dd^nk.
\]

从这里到 \(L^2\) 的延拓用到的是一个标准的完备性论证：若 \(f_m\in\mathcal S\) 在 \(L^2\) 中逼近 \(f\)，则

\[
\|\widehat f_m-\widehat f_n\|_2
=\|f_m-f_n\|_2,
\]

所以 \(\widehat f_m\) 是 \(L^2\) 中的 Cauchy 序列；把其极限定义为 \(\mathcal Ff\)，并由等距性保证定义与逼近序列无关。这样得到的 \(L^2\) Fourier 变换并不要求原积分逐点绝对收敛。

在固定规范后，Fourier 反演、卷积与不确定性关系可以从同一组积分恒等式逐步推出。

前面已经给出了 Fourier 反演与 Plancherel 关系。这里再把两件常被公式化记忆、却很容易混淆归一化常数的事情完整算一遍：第一，为什么反演积分确实回到原函数；第二，位置宽度与波数宽度为什么不能同时任意小。先取 $f\in\mathcal S(\mathbb R^n)$，最后再利用稠密性进入 $L^2$。

直接把 Fourier 定义代回反演式会遇到振荡积分
\[
\int_{\mathbb R^n}\ee^{\ii k\cdot(x-y)}\,\dd^n k,
\]
它不能作为普通 Lebesgue 积分理解。一个稳妥办法是先加入高斯因子，再令正则化参数趋于零。对 $\epsilon>0$ 定义
\[
I_\epsilon(x)
=\frac1{(2\pi)^{n/2}}
\int_{\mathbb R^n}
\widehat f(k)\ee^{\ii k\cdot x}\ee^{-\epsilon|k|^2/2}\,\dd^n k.
\]
把 $\widehat f$ 的定义代入，并利用 Schwartz 衰减交换积分次序，得到
\[
\begin{aligned}
I_\epsilon(x)
&=\frac1{(2\pi)^n}
\int_{\mathbb R^n}\!\int_{\mathbb R^n}
 f(y)\ee^{\ii k\cdot(x-y)}\ee^{-\epsilon|k|^2/2}
\,\dd^n y\,\dd^n k\\
&=\int_{\mathbb R^n}f(y)G_\epsilon(x-y)\,\dd^n y,
\end{aligned}
\]
其中内层高斯积分给出
\begin{equation}
G_\epsilon(z)
=\frac1{(2\pi\epsilon)^{n/2}}
\exp\!\left(-\frac{|z|^2}{2\epsilon}\right).
\label{eq:fourier-gaussian-approx-identity}
\end{equation}
注意
\[
G_\epsilon\ge0,
\qquad
\int_{\mathbb R^n}G_\epsilon(z)\,\dd^n z=1,
\]
并且任意固定 $\delta>0$ 都有
\[
\int_{|z|>\delta}G_\epsilon(z)\,\dd^n z\longrightarrow0,
\qquad \epsilon\downarrow0.
\]
因此 $G_\epsilon$ 是逼近恒等。把积分写成
\[
I_\epsilon(x)-f(x)
=\int G_\epsilon(z)\,[f(x-z)-f(x)]\,\dd^n z,
\]
在 $|z|<\delta$ 与 $|z|\ge\delta$ 两部分分别利用 $f$ 的一致连续性和高斯尾部衰减，就得到
\[
I_\epsilon(x)\longrightarrow f(x).
\]
另一方面，$\ee^{-\epsilon|k|^2/2}\to1$，而 $\widehat f\in\mathcal S\subset L^1$，支配收敛定理又给出
\[
I_\epsilon(x)
\longrightarrow
\frac1{(2\pi)^{n/2}}
\int\widehat f(k)\ee^{\ii k\cdot x}\,\dd^n k.
\]
比较两个极限便得到 Fourier 反演公式。这里的逻辑很重要：所谓
\[
\frac1{(2\pi)^n}\int \ee^{\ii k\cdot(x-y)}\,\dd^n k
=\delta(x-y)
\]
并不是先验的普通积分恒等式，而正可以理解为上述高斯正则化极限的分布写法。

同样的正则化也能把卷积定理中的归一化常数固定下来。本讲义定义
\[
(f*g)(x)
=\frac1{(2\pi)^{n/2}}
\int_{\mathbb R^n}f(y)g(x-y)\,\dd^n y.
\]
于是
\[
\begin{aligned}
\widehat{f*g}(k)
&=\frac1{(2\pi)^{n/2}}
 \int \ee^{-\ii k\cdot x}(f*g)(x)\,\dd^n x\\
&=\frac1{(2\pi)^n}
 \int\!\int \ee^{-\ii k\cdot x}f(y)g(x-y)
 \,\dd^n y\,\dd^n x.
\end{aligned}
\]
令 $z=x-y$，则 $x=z+y$，所以
\[
\begin{aligned}
\widehat{f*g}(k)
&=\left[
\frac1{(2\pi)^{n/2}}
\int f(y)\ee^{-\ii k\cdot y}\,\dd^n y
\right]
\left[
\frac1{(2\pi)^{n/2}}
\int g(z)\ee^{-\ii k\cdot z}\,\dd^n z
\right]\\
&=\widehat f(k)\widehat g(k).
\end{aligned}
\]
反过来，由反演公式可得
\[
\mathcal F^{-1}[\widehat f\,\widehat g]=f*g.
\]
这正是常系数线性方程中``Fourier 空间乘法''与``实空间 Green 卷积''互相转换的根源。

下面把\chapref{chap:hilbert-space}的 Hilbert 空间几何再用一次，推出 Fourier 分析中最重要的不等式之一。先在一维考虑归一化波函数
\[
f\in\mathcal S(\mathbb R),
\qquad \|f\|_2=1.
\]
定义位置与波数的平均值
\[
\bar x=\int x|f(x)|^2\,\dd x,
\qquad
\bar k=\int k|\widehat f(k)|^2\,\dd k,
\]
以及方差
\[
(\Delta x)^2
=\int (x-\bar x)^2|f(x)|^2\,\dd x,
\]
\[
(\Delta k)^2
=\int (k-\bar k)^2|\widehat f(k)|^2\,\dd k.
\]
令
\[
g(x)=\ee^{-\ii\bar k x}f(x).
\]
调制只平移 Fourier 频谱，因此
\[
\|g\|_2=1,
\qquad
\int x|g|^2\,\dd x=\bar x,
\qquad
\|g'\|_2^2=(\Delta k)^2.
\]
最后一个等式由
\[
\mathcal F[g'](k)=\ii k\widehat g(k)
\]
和 Plancherel 定理直接得到。

现在对
\[
(x-\bar x)|g(x)|^2
\]
积分分部。Schwartz 衰减使边界项为零，所以
\[
0
=\int_{-\infty}^{\infty}
\frac{\dd}{\dd x}\big[(x-\bar x)|g|^2\big] \,\dd x.
\]
展开后
\[
0
=1+2\operatorname{Re}
\int (x-\bar x)g^*(x)g'(x)\,\dd x.
\]
因此
\[
\frac12
=\left|
\operatorname{Re}
\int (x-\bar x)g^*g'\,\dd x
\right|
\le
\left|
\int (x-\bar x)g^*g'\,\dd x
\right|.
\]
再用 Cauchy--Schwarz 不等式，得到
\[
\frac12
\le
\|(x-\bar x)g\|_2\,\|g'\|_2
=\Delta x\,\Delta k.
\]
也就是
\begin{equation}
\boxed{\Delta x\,\Delta k\ge\frac12.}
\label{eq:fourier-uncertainty}
\end{equation}
若采用动量 $p=\hbar k$，便成为
\[
\Delta x\,\Delta p\ge\frac\hbar2.
\]

等号什么时候成立也可以从 Cauchy--Schwarz 的等号条件直接读出。必须有
\[
g'(x)=-a(x-\bar x)g(x)
\]
其中 $a>0$ 为实常数。积分得到
\[
g(x)=C\exp\!\left[-\frac a2(x-\bar x)^2\right].
\]
恢复被去掉的平均波数相位后
\[
f(x)=C
\exp\!\left[-\frac a2(x-\bar x)^2+\ii\bar k x\right].
\]
因此高斯函数不是偶然在 Fourier 分析中反复出现：它同时是 Fourier 变换封闭族、热核的基本形状，也是位置--波数不确定关系的极小态。

多维情形对每一个坐标分量完全相同：若
\[
\Delta x_j^2
=\|(x_j-\bar x_j)f\|_2^2,
\qquad
\Delta k_j^2
=\|(k_j-\bar k_j)\widehat f\|_2^2,
\]
则
\[
\Delta x_j\,\Delta k_j\ge\frac12.
\]
更一般地，对任意两个自伴算符 $A,B$，同一套 Hilbert 空间论证会导出 Robertson 不确定关系；这里 Fourier 情形的特殊之处在于 $x$ 与 $-\ii\partial_x$ 恰好通过 Fourier 变换互相对角化。


Plancherel 定理还提供了从 Fourier 空间定义 Sobolev 正则性的最直接方式。对任意实数 $s$，定义 Bessel-potential Sobolev 空间
\begin{equation}
H^s(\RR^n)
:=\left\{u\in\mathcal S'(\RR^n):
(1+|k|^2)^{s/2}\widehat u(k)\in L^2(\RR^n)\right\},
\label{eq:math-dist-sobolev-fourier-definition}
\end{equation}
并赋范
\begin{equation}
\boxed{
\|u\|_{H^s}^2
=\int_{\RR^n}(1+|k|^2)^s|\widehat u(k)|^2\,\dd^n k.
}
\label{eq:math-dist-sobolev-fourier-norm}
\end{equation}
这一定义对所有实数 $s$ 都有意义，因此不仅包含普通函数及其弱导数，也自然包含负阶分布空间。令
\[
\Lambda^s:=(I-\Delta)^{s/2},
\]
则在 Fourier 空间中
\[
\widehat{\Lambda^su}(k)
=(1+|k|^2)^{s/2}\widehat u(k),
\]
所以
\[
\|u\|_{H^s}=\|\Lambda^su\|_{L^2}.
\]
这使 Sobolev 正则性变成一个精确的高频加权：$s$ 越大，对 $|k|\gg1$ 的惩罚越强。

当 $m\in\NN$ 时，Fourier 定义与弱导数定义等价：
\begin{equation}
\boxed{
H^m(\RR^n)
=W^{m,2}(\RR^n)
=\{u\in L^2:D^\alpha u\in L^2,\ |\alpha|\le m\}.
}
\label{eq:math-dist-hm-wm2-equivalence}
\end{equation}
这里 $D^\alpha u$ 按分布导数理解。这个等价性是“广义函数”与“能量空间”真正接上的位置，而不是两个独立定义的巧合。

\begin{derivation}[从 Plancherel 推出 \texorpdfstring{$H^m=W^{m,2}$}{Hm=Wm2}]
由 tempered distribution 的 Fourier 微分律，
\[
\widehat{D^\alpha u}(k)
=(\ii k)^\alpha\widehat u(k).
\]
若 $u\in H^m$，则对任意 $|\alpha|\le m$，由多重指标展开
\[
|k^\alpha|^2
=\prod_{j=1}^n|k_j|^{2\alpha_j}
\le\prod_{j=1}^n(1+k_j^2)^{\alpha_j}
\le(1+|k|^2)^{|\alpha|}
\le(1+|k|^2)^m,
\]
故存在常数 $C_{m,n}$ 使
\[
|k^\alpha|^2
\le C_{m,n}(1+|k|^2)^m.
\]
Plancherel 给出
\[
\|D^\alpha u\|_2^2
=\int|k^\alpha|^2|\widehat u(k)|^2\,\dd^n k
\le C_{m,n}\int(1+|k|^2)^m|\widehat u(k)|^2\,\dd^n k
=C_{m,n}\|u\|_{H^m}^2.
\]
于是所有 $m$ 阶以内分布导数都属于 $L^2$。

反过来，若所有 $|\alpha|\le m$ 的弱导数都在 $L^2$，则由二项式展开
\[
(1+|k|^2)^m
=\sum_{|\alpha|\le m}\binom{m}{\alpha}|k^\alpha|^2
\le C'_{m,n}\sum_{|\alpha|\le m}|k^\alpha|^2,
\]
积分后得到
\[
\|u\|_{H^m}^2
=\int(1+|k|^2)^m|\widehat u(k)|^2\,\dd^n k
\le C'_{m,n}\sum_{|\alpha|\le m}\int|k^\alpha|^2|\widehat u(k)|^2\,\dd^n k
=C'_{m,n}\sum_{|\alpha|\le m}\|D^\alpha u\|_2^2.
\]
两套范数互相控制，因此给出同一个完备空间及等价拓扑。

\emph{举例：$H^1(\RR)$ 与 $W^{1,2}(\RR)$。}对 $u\in H^1(\RR)$，Plancherel 给出
\[
\|u\|_{H^1}^2=\int(1+k^2)|\widehat u(k)|^2\,\dd k
=\|u\|_{L^2}^2+\|u'\|_{L^2}^2,
\]
即 $H^1$ 范数恰是 $L^2$ 范数与导数 $L^2$ 范数的平方和。对 $u(x)=\ee^{-|x|}$，$\widehat u(k)=2/(1+k^2)$，故 $u\in H^1$ 但 $u\notin H^2$（因 $k^2\widehat u(k)\notin L^2$）。

\emph{物理建模。}在量子力学中，$H^1(\RR^3)$ 是氢原子 Hamiltonian 的形式域，波函数 $\psi\in H^1$ 保证动能 $\int|\nabla\psi|^2$ 有限。Sobolev 嵌入 $H^1(\RR^3)\hookrightarrow L^6(\RR^3)$ 则给出 Hartree 非线性项 $\int|\psi|^4$ 的良定义性，这是多体 Schrödinger 方程能谱分析的基础。
\end{derivation}

Fourier 权重还直接给出最基本的 Sobolev embedding。若
\[
s>\frac n2,
\]
则 Cauchy--Schwarz 产生
\begin{align*}
\int_{\RR^n}|\widehat u(k)|\,\dd^n k
&\le
\left[
\int(1+|k|^2)^{-s}\,\dd^n k
\right]^{1/2}
\|u\|_{H^s}.
\end{align*}
括号中的积分恰在 $s>n/2$ 时有限。Fourier 反演于是给出
\begin{equation}
\boxed{
H^s(\RR^n)\hookrightarrow C_b^0(\RR^n),
\qquad s>\frac n2.
}
\label{eq:math-dist-sobolev-c0-embedding}
\end{equation}
更一般地，若
\[
s>k+\frac n2,
\]
则
\begin{equation}
H^s(\RR^n)\hookrightarrow C_b^k(\RR^n).
\label{eq:math-dist-sobolev-ck-embedding}
\end{equation}
因为 $D^\alpha u$ 的 Fourier transform 额外带 $k^\alpha$，只需把上面的可积权重改成
\[
|k|^{2|\alpha|}(1+|k|^2)^{-s}.
\]
这个阈值 $s>n/2$ 也是非线性 PDE 中 Sobolev 空间开始成为乘法代数的临界来源：足够多的平方可积导数最终控制点态乘法。

常系数椭圆算子的正则性也可以先在 Fourier 空间看清。最透明的母模型是
\[
I-\Delta.
\]
若
\begin{equation}
(I-\Delta)u=f,
\label{eq:math-dist-bessel-elliptic-equation}
\end{equation}
则
\[
(1+|k|^2)\widehat u(k)=\widehat f(k),
\]
从而精确地有
\begin{equation}
\boxed{
\|u\|_{H^{s+2}}
=\|f\|_{H^s}.
}
\label{eq:math-dist-bessel-two-derivative-gain}
\end{equation}
也就是说，椭圆逆算子在 Sobolev 标度上提升两个导数。对一般阶数 $m$ 的常系数椭圆算子 $P(D)$，若主符号满足
\[
|p(k)|\ge c|k|^m
\qquad(|k|\gg1),
\]
则同样的高频分解给出先验估计
\begin{equation}
\|u\|_{H^{s+m}}
\le C\bigl(\|P(D)u\|_{H^s}+\|u\|_{H^s}\bigr).
\label{eq:math-dist-constant-elliptic-estimate}
\end{equation}
低频处附加的 $\|u\|_{H^s}$ 正是可能存在 kernel 或小频退化时必须保留的项。变系数椭圆理论把这条 Fourier 估计局部化，并用 commutator 与 partition of unity 控制误差；\chapref{chap:nonlinear-pde} 中的能量估计正是在这一分析接口之上继续发展。

接下来考察卷积。

本讲义取

\[
(f*g)(x)=\frac1{(2\pi)^{n/2}}\int_{\mathbb R^n}f(y)g(x-y)\,\dd^n y.
\]

在这个归一化下

\[
\mathcal F[f*g]=\widehat f\,\widehat g.
\]

对分布，卷积并非总有定义。一个常用充分条件是：两个分布中至少有一个具有紧支集。此时卷积可通过张量积分严格定义。

\section{奇异核、基本解与逼近恒等}

前面的定义已经能处理 \(\delta\)、主值分布和 Fourier 变换。真正进入数学物理时，还需要看清三件事：奇异边界值怎样产生 \(\delta\) 项；微分算符的基本解怎样由 Fourier 空间构造；Green 函数为什么可以看成“带边界条件的基本解”。

接下来考察奇异核、Heaviside 与主值分布。

Heaviside 分布 \(H(x)\) 满足
\[
H'(x)=\delta(x).
\]
符号函数
\[
\operatorname{sgn}x=2H(x)-1
\]
因此在分布意义下
\[
\boxed{
\frac{\dd}{\dd x}\operatorname{sgn}x=2\delta(x)
}.
\]
这个等式说明：经典函数在跳跃点的“缺失导数”会以 \(\delta\) 的形式重新进入分布导数。

再看
\[
\ln|x|.
\]
它局部可积，因此定义正则分布。对检验函数 \(\phi\)，
\[
\left\langle \frac{\dd}{\dd x}\ln|x|,\phi\right\rangle
=-\int_{-\infty}^{\infty}\ln|x|\,\phi'(x)\,\dd x.
\]
把积分在 \(( -\infty,-\varepsilon)\cup(\varepsilon,\infty)\) 上分部积分，得到
\[
\begin{aligned}
&-\int_{|x|>\varepsilon}\ln|x|\,\phi'(x)\,\dd x\\
&\quad=\int_{|x|>\varepsilon}\frac{\phi(x)}{x}\,\dd x
+\ln\varepsilon\,[\phi(-\varepsilon)-\phi(\varepsilon)].
\end{aligned}
\]
因为
\[
\phi(-\varepsilon)-\phi(\varepsilon)=O(\varepsilon),
\]
第二项在 \(\varepsilon\downarrow0\) 时趋于零。故
\begin{equation}
\boxed{
\frac{\dd}{\dd x}\ln|x|
=\operatorname{pv}\frac1x
}.
\label{eq:log-pv-derivative}
\end{equation}
这里绝不能把右边当成普通函数 \(1/x\) 的 Lebesgue 积分；主值极限是定义的一部分。

Sokhotski--Plemelj 公式进一步把主值部分与 delta 奇性统一到同一个正则化极限中。

考虑
\[
\frac1{x\pm \ii\varepsilon}
=\frac{x}{x^2+\varepsilon^2}
\mp \ii\frac{\varepsilon}{x^2+\varepsilon^2}.
\]
第二项包含熟悉的 Poisson 核。对任意检验函数 \(\phi\)，令 \(x=\varepsilon y\)，有
\[
\begin{aligned}
\int_{-\infty}^{\infty}
\frac{\varepsilon}{x^2+\varepsilon^2}\phi(x)\,\dd x
&=\int_{-\infty}^{\infty}\frac{\phi(\varepsilon y)}{1+y^2}\,\dd y\\
&\longrightarrow \phi(0)\int_{-\infty}^{\infty}\frac{\dd y}{1+y^2}\\
&=\pi\phi(0).
\end{aligned}
\]
因此
\[
\frac{\varepsilon}{x^2+\varepsilon^2}
\longrightarrow \pi\delta(x)
\]
于 \(\mathcal D'\)。

实部则趋于 Cauchy 主值：
\[
\frac{x}{x^2+\varepsilon^2}
\longrightarrow \operatorname{pv}\frac1x.
\]
于是得到边界值公式
\begin{equation}
\boxed{
\frac1{x\pm \ii 0}
=\operatorname{pv}\frac1x\mp \ii\pi\delta(x)
}.
\label{eq:sokhotski-plemelj}
\end{equation}
这不是记号技巧，而是复平面上从上、下半平面逼近实轴时，两种边界值分布的精确差别。

平移立即给出
\[
\frac1{x-x_0\pm \ii 0}
=\operatorname{pv}\frac1{x-x_0}\mp \ii\pi\delta(x-x_0).
\]
后面处理 Green 函数的因果边界条件、Hilbert 变换以及频域积分时，这个公式会反复出现。

接下来考察Fourier 方法与微分算符基本解。

设常系数微分算符
\[
P(D)=\sum_{|\alpha|\le m}a_\alpha D^\alpha,
\]
希望寻找分布 \(E\) 使
\[
P(D)E=\delta.
\]
对两边作 Fourier 变换，得到
\[
P(\ii k)\widehat E(k)=\frac1{(2\pi)^{n/2}}.
\]
形式上
\begin{equation}
\boxed{
\widehat E(k)
=\frac1{(2\pi)^{n/2}}\frac1{P(\ii k)}
}.
\label{eq:fundamental-solution-fourier}
\end{equation}
真正的困难集中到了符号 \(P(\ii k)\) 的零点：若分母从不为零，直接逆变换即可；若它在某些 \(k\) 上为零，就必须指定主值、\(\pm \ii 0\) 或其他边界值处方。不同处方对应不同的物理解，例如 retarded、advanced 或 outgoing Green 函数。

以一维 Helmholtz 型算符为例，考虑
\[
\left(\frac{\dd^2}{\dd x^2}+k_0^2\right)G(x)=\delta(x),
\qquad k_0>0.
\]
在 \(x\neq0\) 区域，齐次解是 \(\ee^{\pm \ii k_0x}\)。若要求向外传播，可取
\[
G_+(x)=C \ee^{\ii k_0|x|}.
\]
连续性保证 \(G_+\) 在 \(0\) 处不跳跃。把方程在 \((-\varepsilon,\varepsilon)\) 上积分：
\[
G_+'(\varepsilon)-G_+'(-\varepsilon)
+k_0^2\int_{-\varepsilon}^{\varepsilon}G_+(x)\,\dd x=1.
\]
令 \(\varepsilon\downarrow0\)，第二项消失，故
\[
G_+'(0^+)-G_+'(0^-)=1.
\]
而
\[
G_+'(0^+)=\ii k_0C,
\qquad
G_+'(0^-)=-\ii k_0C,
\]
所以
\[
2\ii k_0C=1,
\qquad
C=\frac1{2\ii k_0}.
\]
因此
\begin{equation}
\boxed{
G_+(x)=\frac{\ee^{\ii k_0|x|}}{2\ii k_0}
}.
\label{eq:1d-outgoing-helmholtz-green}
\end{equation}
这个“导数跳跃 = 点源强度”的推导比单纯查表更重要，因为它可以直接推广到分段 Green 函数的构造。

三维 Poisson 方程给出这一一般 Fourier 基本解方法的标准椭圆型特例。

在三维中考虑
\[
-\nabla^2G(\mathbf x)=\delta(\mathbf x).
\]
远离原点时 \(G\) 调和，而且旋转对称性提示
\[
G(\mathbf x)=g(r),
\qquad r=|\mathbf x|.
\]
对 \(r>0\)，
\[
\nabla^2g
=\frac1{r^2}\frac{\dd}{\dd r}
\left(r^2\frac{\dd g}{\dd r}\right)=0.
\]
积分两次得到
\[
g(r)=A+\frac{B}{r}.
\]
若要求无穷远衰减，则 \(A=0\)。系数由点源归一化确定。对半径 \(\varepsilon\) 的球积分，使用 Gauss 定理：
\[
1=\int_{B_\varepsilon}\delta(\mathbf x)\,\dd^3x
=-\int_{B_\varepsilon}\nabla^2G\,\dd^3x
=-\oint_{S_\varepsilon}\nabla G\cdot\dd\mathbf S.
\]
若 \(G=B/r\)，则
\[
\nabla G=-\frac{B}{r^2}\hat{\mathbf r},
\]
所以
\[
1=-\left(-\frac{B}{\varepsilon^2}\right)4\pi\varepsilon^2
=4\pi B.
\]
因此
\begin{equation}
\boxed{
G(\mathbf x)=\frac1{4\pi|\mathbf x|},
\qquad
-\nabla^2G=\delta
}.
\label{eq:poisson-fundamental-solution}
\end{equation}
于是全空间 Poisson 方程
\[
-\nabla^2u=f
\]
的形式精确解为
\[
u(\mathbf x)=\int_{\mathbb R^3}
\frac{f(\mathbf y)}{4\pi|\mathbf x-\mathbf y|}\,\dd^3y,
\]
只要函数空间和无穷远条件使卷积有意义。

接下来考察热核作为时空分布的基本解。

热方程
\[
(\partial_t-\kappa\nabla^2)u=0
\]
的 Cauchy 问题可以从 Fourier 空间完全解出。对空间变量作 Fourier 变换：
\[
\partial_t\widehat u(k,t)+\kappa|k|^2\widehat u(k,t)=0.
\]
这是关于 \(t\) 的一阶常微分方程，故
\[
\widehat u(k,t)=\ee^{-\kappa|k|^2t}\widehat u_0(k).
\]
逆变换写成卷积
\[
u(x,t)=G_t*u_0,
\]
其中热核为
\begin{equation}
\boxed{
G_t(x)=\frac1{(4\pi\kappa t)^{n/2}}
\exp\!\left(-\frac{|x|^2}{4\kappa t}\right),
\qquad t>0
}.
\label{eq:heat-kernel}
\end{equation}
并满足
\[
\int_{\mathbb R^n}G_t(x)\,\dd^nx=1.
\]
当 \(t\downarrow0\) 时，\(G_t\) 不逐点趋于一个普通函数，但在分布意义下
\[
G_t\longrightarrow\delta.
\]
所以更精确地说，因果热 Green 函数
\[
G(x,t)=H(t)G_t(x)
\]
满足
\[
(\partial_t-\kappa\nabla^2)G
=\delta(t)\delta(x).
\]
这把“初值传播核”与“微分算符基本解”统一成了同一件事。

接下来考察Green 函数与基本解的区别。

基本解只要求
\[
L E=\delta
\]
在全空间分布意义下成立；Green 函数还要满足具体区域的边界条件。因此同一个微分算符可以有同一个局部奇异结构，却因边界条件不同得到不同 Green 函数。

设 \(\Omega\) 是有边界区域，\(G(x,y)\) 满足
\[
L_xG(x,y)=\delta(x-y)
\]
以及指定边界条件。若 \(L\) 对应的边值问题可逆，则
\[
u(x)=\int_\Omega G(x,y)f(y)\,\dd y
\]
是 \(Lu=f\) 的解。若 \(L\) 有零模，则这个逆不再存在于整个空间上，此时必须像 Fredholm alternative 那样加入与伴随零模正交的可解条件。\chapref{chap:integral-equations}会从积分算符角度重新得到同一结论。

从全空间常系数基本解继续向变系数问题推进时，普通 Fourier multiplier 已经不够。真正需要的是“允许 multiplier 同时依赖位置”的伪微分算子。它把前面的 Sobolev 权重、wavefront set 与椭圆正则性统一到同一套 symbol calculus 中。

设
\[
\langle\xi\rangle=(1+|\xi|^2)^{1/2}.
\]
对实数 $m$，Kohn--Nirenberg 符号类 $S^m_{1,0}(\RR^n_x\times\RR^n_\xi)$ 由所有光滑函数 $a(x,\xi)$ 组成，并满足对任意 multi-indices $\alpha,\beta$ 都存在 $C_{\alpha\beta}$ 使
\begin{equation}
\boxed{
|\partial_x^\alpha\partial_\xi^\beta a(x,\xi)|
\le C_{\alpha\beta}\langle\xi\rangle^{m-|\beta|}.
}
\label{eq:math-dist-symbol-class}
\end{equation}
$x$-导数不降低阶数，而每作一次 $\xi$-导数都降低一个高频阶数。这正是“局部系数变化不改变微分阶，而频率微分会削弱高频增长”的量化版本。

给定 $a\in S^m_{1,0}$，定义伪微分算子
\begin{equation}
\boxed{
\operatorname{Op}(a)u(x)
=\frac1{(2\pi)^n}
\iint
\ee^{\ii(x-y)\cdot\xi}
a(x,\xi)u(y)
\,\dd^ny\,\dd^n\xi.
}
\label{eq:math-dist-pseudodifferential-quantization}
\end{equation}
对 Schwartz 函数该式可先按振荡积分理解。若 $a$ 与 $x$ 无关，就退化为 Fourier multiplier；若
\[
a(x,\xi)=\sum_{|\alpha|\le m}a_\alpha(x)(\ii\xi)^\alpha,
\]
则 $\operatorname{Op}(a)$ 正是普通微分算子
\[
P(x,D)=\sum_{|\alpha|\le m}a_\alpha(x)D^\alpha.
\]
因此 differential operator 只是 symbol 在 $\xi$ 上为多项式的特殊伪微分算子。

若两个 $m$ 阶符号之差属于 $S^{m-1}$，它们具有相同 principal symbol。于是 principal symbol 应理解为商类
\begin{equation}
\sigma_m(P)\in S^m/S^{m-1},
\label{eq:math-dist-principal-symbol-quotient}
\end{equation}
而不是依赖某个低阶写法的具体函数。对微分算子，上式正恢复最高阶齐次多项式
\[
p_m(x,\xi)
=\sum_{|\alpha|=m}a_\alpha(x)(\ii\xi)^\alpha.
\]
高频极限只看 principal symbol，这正是椭圆性、特征集与奇性传播全部由最高阶部分控制的原因。

伪微分算子的第一条基本映射性质是
\begin{equation}
\boxed{
P\in\Psi^m
\quad\Longrightarrow\quad
P:H^s_{\mathrm{comp}}
\longrightarrow H^{s-m}_{\mathrm{loc}}
}
\label{eq:math-dist-pseudo-sobolev-mapping}
\end{equation}
在 proper-support 或适当局部化条件下连续。$m>0$ 的算子损失 $m$ 阶 Sobolev 正则性，$m<0$ 的算子则获得 $|m|$ 阶正则性；$S^{-\infty}:=\cap_mS^m$ 对应 smoothing operators，把任意 distribution 局部送到 $C^\infty$。

更重要的是 symbol calculus 对复合封闭。若
\[
A=\operatorname{Op}(a),
\qquad
B=\operatorname{Op}(b),
\]
则 $AB=\operatorname{Op}(a\# b)$，其中有渐近展开
\begin{equation}
\boxed{
a\# b
\sim
\sum_{\alpha}
\frac{1}{\alpha!}\left(\frac1\ii\right)^{|\alpha|}
(\partial_\xi^\alpha a)
(\partial_x^\alpha b).
}
\label{eq:math-dist-symbol-composition}
\end{equation}
“渐近”表示截断到 $|\alpha|<N$ 后的余项降低 $N$ 个 symbol 阶。零阶项只是 $ab$，因此
\[
\sigma_{m_1+m_2}(AB)
=\sigma_{m_1}(A)\sigma_{m_2}(B).
\]

\begin{derivation}[伪微分复合为何产生 \texorpdfstring{$\partial_\xi a\,\partial_xb$}{frequency-position derivatives}]
\emph{策略：}把复合算子写成双重积分核，做一次变量替换分离出一个「快振荡相位」，再把符号沿位置 Taylor 展开，并用振荡相位把位置多项式转成频率导数。

\emph{第一步：双重积分核。}把两个 kernel 相乘，得到
\begin{align*}
ABu(x)
={}&\frac1{(2\pi)^{2n}}
\iiint\!\int
\ee^{\ii(x-y)\cdot\xi}
a(x,\xi)
\ee^{\ii(y-z)\cdot\eta}
b(y,\eta)u(z)
\,\dd z\,\dd\eta\,\dd y\,\dd\xi.
\end{align*}

\emph{第二步：分离快相位。}令 $\theta=\xi-\eta$，即 $\xi=\eta+\theta$。总相位变为
\[
(x-y)\cdot(\eta+\theta)+(y-z)\cdot\eta
=(x-z)\cdot\eta+(x-y)\cdot\theta.
\]
这一步的意图是：$z$ 只出现在 $(x-z)\cdot\eta$ 里，$y$ 只出现在 $(x-y)\cdot\theta$ 里，两者解耦，为后续分别对 $z$、$y$ 积分做准备。

\emph{第三步：沿位置 Taylor 展开。}因为 $\ee^{\ii(x-y)\cdot\theta}$ 对 $y$ 是快振荡的，$y$ 远离 $x$ 时正负相位相消（稳定相位原理），主导贡献来自 $y\approx x$。故把 $b(y,\eta)$ 在 $y=x$ 处展开：
\[
b(y,\eta)
\sim\sum_\alpha
\frac{(y-x)^\alpha}{\alpha!}
\partial_x^\alpha b(x,\eta).
\]

\emph{第四步：多项式换成频率导数。}位置多项式乘上振荡指数可改写为对 $\theta$ 的导数：
\[
(y-x)^\alpha\ee^{\ii(x-y)\cdot\theta}
=\left(\frac1\ii\right)^{|\alpha|}
\partial_\theta^\alpha
\ee^{\ii(x-y)\cdot\theta},
\]
因为 $\partial_\theta$ 每作用一次提出一个 $\ii(x-y)$ 因子。

\emph{第五步：分部积分移位。}把 $\partial_\theta^\alpha$ 从指数上分部积分移到 $a(x,\eta+\theta)$ 上（边界项因符号速降而消失，并注意每次分部积分产生一个负号）：
\begin{equation*}
 \int a(x,\eta+\theta)\,\partial_\theta^\alpha \ee^{\ii(x-y)\cdot\theta}\,\dd\theta
 =(-1)^{|\alpha|}
 \int \partial_\theta^\alpha a(x,\eta+\theta)\,\ee^{\ii(x-y)\cdot\theta}\,\dd\theta.
\end{equation*}

\emph{第六步：对 $y$、$z$ 积分并置 $\theta=0$。}对 $z$ 积分给出 $\hat u(\eta)$，对 $y$ 积分给出 $(2\pi)^n\delta(\theta)$，于是令 $\theta=0$、\(\partial_\theta^\alpha a(x,\eta+\theta)|_{\theta=0}=\partial_\xi^\alpha a(x,\eta)\)，得到
\[
\sum_\alpha
\frac1{\alpha!}
\left(\frac1\ii\right)^{|\alpha|}
\partial_\xi^\alpha a(x,\eta)
\partial_x^\alpha b(x,\eta),
\]
即\eqnrefs{eq:math-dist-symbol-composition}。所以 symbol calculus 本质上是“振荡相位把位置 Taylor 系数转成频率导数”。

\emph{一阶项与 Poisson 括号。}取 $|\alpha|=1$，一阶修正为
\[
\frac1\ii\sum_j \partial_{\xi_j}a\,\partial_{x^j}b,
\]
即算子复合的领头量子修正由「频率方向导数 × 位置方向导数」给出。这正是下一段交换子中 Poisson 括号\eqnrefs{eq:math-dist-commutator-poisson} 的来源：零阶项 $ab$ 在交换子 $AB-BA$ 中相消，剩下一阶项给出 $\frac1\ii\{a,b\}$。
\end{derivation}

交换子的最高阶会自动消去，因为 $ab-ba=0$。若 $A\in\Psi^{m_1}$、$B\in\Psi^{m_2}$ 为标量算子，则
\begin{equation}
[A,B]\in\Psi^{m_1+m_2-1},
\label{eq:math-dist-pseudo-commutator-order}
\end{equation}
而其 principal symbol 为
\begin{equation}
\boxed{
\sigma_{m_1+m_2-1}([A,B])
=\frac1\ii\{a_{m_1},b_{m_2}\},
}
\label{eq:math-dist-commutator-poisson}
\end{equation}
其中
\[
\{a,b\}
=\sum_j
\left(
\partial_{\xi_j}a\,\partial_{x^j}b
-\partial_{x^j}a\,\partial_{\xi_j}b
\right)
\]
是余切丛上的 canonical Poisson bracket。于是“量子交换子降低一阶并趋向 classical Poisson bracket”并不是半经典理论里突然出现的类比，而已经内嵌在 ordinary pseudodifferential symbol calculus 中。

现在回到正则性。若 $P\in\Psi^m$ 的 principal symbol $p_m(x,\xi)$ 在某个 conic 区域满足
\begin{equation}
|p_m(x,\xi)|\ge c|\xi|^m
\qquad(|\xi|\gg1),
\label{eq:math-dist-microlocal-ellipticity}
\end{equation}
就称 $P$ 在该区域 microlocally elliptic。若处处满足，则称 elliptic。其 characteristic set 为
\begin{equation}
\operatorname{Char}(P)
:=\{(x,\xi)\in T^*M\setminus0:p_m(x,\xi)=0\}.
\label{eq:math-dist-characteristic-set}
\end{equation}
椭圆算子正是 $\operatorname{Char}(P)=\varnothing$ 的情形。

椭圆性的核心不是一个不等式本身，而是可以逐阶构造 parametrix。先取
\[
q_{-m}^{(0)}(x,\xi)
=\frac{\chi(x,\xi)}{p_m(x,\xi)},
\]
其中 $\chi$ 在目标 conic 区域的大频率处等于 $1$。则
\[
q_{-m}^{(0)}\# p=1+r_{-1},
\qquad r_{-1}\in S^{-1}.
\]
继续递归选择 $q_{-m-j}$ 消去 $r$ 的每一阶，可得到 $q\in S^{-m}$ 使
\begin{equation}
\boxed{
QP=I-R,
\qquad
PQ=I-S,
\qquad
R,S\in\Psi^{-\infty}.
}
\label{eq:math-dist-elliptic-parametrix}
\end{equation}
这就是 parametrix：它未必是全局精确逆，但误差已经 smooth 到对 wavefront set 完全不可见。

\begin{derivation}[parametrix 怎样推出椭圆正则性]
\emph{策略：}先用 parametrix 恒等式把解 \(u\) 分解为「正则部分 \(Qf\)」与「光滑误差 \(Ru\)」，再借 \(Q\) 的负阶符号提升 Sobolev 指数，最后用 bootstrap 归纳与 wavefront 论证闭合正则性。

\emph{第一步：解的分裂。}设 $Pu=f$。把 $QP=I-R$ 左作用到方程上，得
\begin{equation*}
  QPu=QP u = (I-R)u = Qf,
\end{equation*}
移项即
\begin{equation*}
  u=Qf+Ru.
\end{equation*}
这一步把未知解 \(u\) 拆成两项：\(Qf\) 由已知数据 \(f\) 经 parametrix \(Q\) 决定，\(Ru\) 是「光滑到不可见」的误差项。

\emph{第二步：\(Q\) 的 Sobolev 提升。}\(Q\in\Psi^{-m}\) 的符号 \(q\in S^{-m}\) 满足
\begin{equation*}
  |q(x,\xi)|\le C(1+|\xi|)^{-m},
\end{equation*}
而 Sobolev 空间 \(H^s\) 的定义是 \((1+|\xi|)^s\hat u\in L^2\)。于是 \(\widehat{Q u}=q\hat u\)（主部）满足
\begin{equation*}
  (1+|\xi|)^{s+m}|\widehat{Q u}|
  \le C(1+|\xi|)^{s}|\hat u|
  \in L^2,
\end{equation*}
即 \(Q u\in H^{s+m}\)。直观上，每个 \(m\) 阶负号相当于「抹平」了 \(m\) 阶导数，因此 parametrix 把正则指数整体抬高 \(m\)。误差算子 \(R\in\Psi^{-\infty}\) 则把任意分布映成光滑函数，不损失任何正则性。故若 \(f\in H^s_{\mathrm{loc}}\)，则
\begin{equation}
 u=Qf+Ru\in H^{s+m}_{\mathrm{loc}},
 \label{eq:math-dist-elliptic-sobolev-regularity}
\end{equation}
此即在 elliptic 区域成立的「正则性提升」。

\emph{第三步：bootstrap 归纳。}若 \(f\) 光滑，先取 \(s\) 足够小使 \(u\in H^s\)（任何分布局部地属于某个 \(H^s\)），第一步给 \(u\in H^{s+m}\)；把新的 \(u\in H^{s+m}\) 代回原方程 \(Pu=f\) 再应用第一步，得 \(u\in H^{s+2m}\)；依此类推，对任意 \(k\) 有 \(u\in H^{s+km}\)。由 Sobolev 嵌入 \(H^\infty\subset C^\infty\)，最终 \(u\) 光滑。

\emph{第四步：wavefront set 的精确化。}从 \(u=Qf+Ru\)，因 \(Ru\) 光滑，其波前集为空，故
\begin{equation*}
  \operatorname{WF}(u)
  =\operatorname{WF}(Qf)
  \subseteq\operatorname{WF}(f)\cup\operatorname{Char}(P),
\end{equation*}
其中右端补上 \(\operatorname{Char}(P)\) 是因为 \(Q\) 只在 \(\operatorname{Char}(P)\) 之外才是真正的逆；在特征集上 \(Q\) 不再控制 \(P\) 的方向。另一方面，任何 pseudodifferential operator 都不会创造原来不存在的波前方向：
\[
 \operatorname{WF}(Pu)\subseteq\operatorname{WF}(u).
\]
在 \(\operatorname{Char}(P)\) 之外，第一条给出
\[
 \operatorname{WF}(u)\subseteq\operatorname{WF}(Pu),
\]
与第二条合并，即在 elliptic 区域
\begin{equation}
 \boxed{
 \operatorname{WF}(u)=\operatorname{WF}(Pu).
 }
 \label{eq:math-dist-wavefront-elliptic-regularity}
\end{equation}
这正把本章前面的 wavefront 定义闭合成一个可计算的 PDE 定理：椭圆区域不改变解的奇异方向。

\emph{举例：Laplace 算子的经典椭圆正则性。}取 \(P=-\Delta\)（\(m=2\)），它是处处椭圆算子，\(\operatorname{Char}(-\Delta)=\varnothing\)。于是 \(\operatorname{WF}(u)=\operatorname{WF}(\Delta u)\)：若 \(\Delta u=f\) 在某区域光滑，则 \(u\) 在该区域光滑。这正是位势论中「调和函数必解析、Poisson 方程右端光滑则解光滑」的严格微局域表述；对一般 \(m\) 阶椭圆算子，\(u\) 比 \(f\) 多 \(m\) 阶导数，且奇异方向在椭圆区域精确保持。
\end{derivation}

当 $P$ 不椭圆时，真正困难只剩在 $\operatorname{Char}(P)$ 上。若 principal symbol $p$ 为实值且满足 real-principal-type 的非退化条件，则 Hamilton vector field
\begin{equation}
H_p
=\sum_j
\left(
\frac{\partial p}{\partial\xi_j}\frac{\partial}{\partial x^j}
-\frac{\partial p}{\partial x^j}\frac{\partial}{\partial\xi_j}
\right)
\label{eq:math-dist-hamilton-symbol-flow}
\end{equation}
在 characteristic set 上生成 bicharacteristics。propagation of singularities theorem 的核心结论是：若 $Pu=f$，那么在 $\operatorname{WF}(f)$ 之外，$\operatorname{WF}(u)\cap\operatorname{Char}(P)$ 必沿 $H_p$ 的 integral curves 传播。椭圆算子没有 characteristic directions，所以奇性被完全压制；双曲算子则允许奇性沿特征流传播。

以波算子
\[
\Box=\partial_t^2-c^2\Delta_x
\]
为例，principal symbol 可取
\[
p(t,x;\tau,\xi)
=-\tau^2+c^2|\xi|^2.
\]
characteristic set 满足
\[
\tau^2=c^2|\xi|^2,
\]
正对应有限传播速度的光锥方向。于是“波前沿光线传播”在微局部语言中就是 Hamiltonian flow of the principal symbol；这同时解释了几何光学、经典 Hamilton flow 与双曲 PDE 奇性传播为何共享同一余切丛结构。

至此，本章的逻辑闭环变成
\[
\begin{aligned}
\text{distribution}
&\to\text{Fourier/Sobolev regularity}
\to\text{wavefront set}
\to\text{symbol calculus},\\
&\to\text{parametrix}
\to\text{elliptic regularity / characteristic propagation}.
\end{aligned}
\]
后续 PDE 章节可以直接调用这套母结构，而不必把“弱导数、频率衰减、Green 函数和正则性”当作互不相干的技巧。


高斯核提供了一族完全由普通函数构成的逼近恒等，可用来检验前述分布极限。

高斯函数在 Fourier 分析中有一个特殊地位：它在 Fourier 变换下仍然是高斯，而且经过尺度压缩后会在分布意义下逼近 \(\delta\)。这两点把 Plancherel、热核和分布极限连在一起。

先从一维高斯的 Fourier 变换开始。

先计算
\[
I(k)=\int_{-\infty}^{\infty}\ee^{-a x^2}\ee^{-\ii kx}\,\dd x,
\qquad a>0.
\]
不直接配复平方也可以通过微分方程求出。对 \(k\) 求导：
\[
I'(k)
=-\ii\int x \ee^{-a x^2}\ee^{-\ii kx}\,\dd x.
\]
利用
\[
x \ee^{-a x^2}
=-\frac1{2a}\frac{\dd}{\dd x}\ee^{-a x^2},
\]
得到
\[
I'(k)
=\frac{\ii}{2a}
\int\frac{\dd}{\dd x}\ee^{-a x^2}\ee^{-\ii kx}\,\dd x.
\]
分部积分，边界项由高斯衰减消失：
\[
I'(k)
=-\frac{k}{2a}I(k).
\]
所以
\[
I(k)=C \ee^{-k^2/(4a)}.
\]
在 \(k=0\) 处用标准高斯积分
\[
I(0)=\sqrt{\frac\pi a}
\]
确定常数，得到
\[
\boxed{
\int_{-\infty}^{\infty}\ee^{-a x^2}\ee^{-\ii kx}\,\dd x
=\sqrt{\frac\pi a}
\ee^{-k^2/(4a)}
}.
\]
在本讲义对称 Fourier 归一化下，
\begin{equation}
\boxed{
\mathcal F[\ee^{-a x^2}](k)
=\frac1{\sqrt{2a}}\ee^{-k^2/(4a)}
}.
\label{eq:gaussian-fourier}
\end{equation}
多维情形因各坐标积分分离而直接得到
\[
\mathcal F[\ee^{-a|x|^2}](k)
=(2a)^{-n/2}\ee^{-|k|^2/(4a)}.
\]

缩放高斯族在分布意义下收敛到 delta。

定义
\[
\rho_\varepsilon(x)
=\frac1{\sqrt{2\pi}\varepsilon}
\exp\!\left(-\frac{x^2}{2\varepsilon^2}\right).
\]
它满足
\[
\rho_\varepsilon(x)\ge0,
\qquad
\int\rho_\varepsilon(x)\,\dd x=1.
\]
对检验函数 \(\phi\)，令 \(x=\varepsilon y\)，得到
\[
\begin{aligned}
\langle\rho_\varepsilon,\phi\rangle
&=\int_{-\infty}^{\infty}
\frac{\ee^{-y^2/2}}{\sqrt{2\pi}}
\phi(\varepsilon y)\,\dd y.
\end{aligned}
\]
当 \(\varepsilon\downarrow0\) 时，\(\phi(\varepsilon y)\to\phi(0)\)。又因为 \(\phi\) 有界，高斯权重可积，可用支配收敛定理得到
\[
\lim_{\varepsilon\downarrow0}
\langle\rho_\varepsilon,\phi\rangle
=\phi(0).
\]
因此
\begin{equation}
\boxed{
\rho_\varepsilon\longrightarrow\delta
\quad\text{于 }\mathcal D'(\mathbb R)
}.
\label{eq:gaussian-delta-sequence}
\end{equation}

注意这个极限不是逐点极限。对任何固定 \(x\neq0\)，
\[
\rho_\varepsilon(x)\to0,
\]
而在 \(x=0\) 处
\[
\rho_\varepsilon(0)\to\infty.
\]
真正稳定不变的是总积分为 1，以及对测试函数的作用趋于取值泛函。

更一般地，这一机制就是逼近恒等与卷积恢复。

对适当函数 \(f\)，定义
\[
f_\varepsilon=\rho_\varepsilon*f.
\]
若采用普通卷积归一化写法，
\[
f_\varepsilon(x)
=\int\rho_\varepsilon(x-y)f(y)\,\dd y.
\]
这可以理解为在宽度 \(\varepsilon\) 上对 \(f\) 做局部平均。当 \(f\) 连续时，固定 \(x\) 有
\[
f_\varepsilon(x)\to f(x).
\]
若 \(f\in L^p\)、\(1\le p<\infty\)，则在标准逼近恒等条件下还有
\[
\|f_\varepsilon-f\|_p\to0.
\]

Fourier 空间中，这个过程更简单。由 \cref{eq:gaussian-fourier}，高斯卷积相当于乘上一个随 \(\varepsilon\to0\) 趋于 1 的高频滤波因子。因此“卷积平滑”与“Fourier 高频抑制”是同一个过程的两种表述。

对上述逼近恒等求导，还可以得到 delta 各阶导数的普通函数逼近。

既然
\[
\rho_\varepsilon\to\delta,
\]
分布微分的连续性给出
\[
\rho_\varepsilon^{(m)}\to\delta^{(m)}.
\]
直接验证也很简单：
\[
\langle\rho_\varepsilon',\phi\rangle
=-\langle\rho_\varepsilon,\phi'\rangle
\longrightarrow-\phi'(0)
=\langle\delta',\phi\rangle.
\]
因此 \(\delta'\)、\(\delta''\) 等奇异分布也可以看成一族越来越尖锐、具有正负振荡结构的普通光滑函数极限。


接下来考察群同态、核与第一同构定理。

前面讨论群同构时，真正更一般的对象是群同态。映射
\[
\varphi:G\to H
\]
若满足
\[
\varphi(g_1g_2)=\varphi(g_1)\varphi(g_2),
\]
就称为群同态。定义
\[
\ker\varphi=
\{g\in G:\varphi(g)=e_H\},
\]
\[
\operatorname{Im}\varphi
=\{\varphi(g):g\in G\}.
\]
像构成 \(H\) 的子群，而核不仅是子群，还是正规子群。确实，若 \(k\in\ker\varphi\)、\(g\in G\)，则
\[
\varphi(gkg^{-1})
=\varphi(g)e_H\varphi(g)^{-1}=e_H,
\]
故
\[
gkg^{-1}\in\ker\varphi.
\]
因此
\[
\boxed{\ker\varphi\trianglelefteq G}.
\]

同态把“哪些元素在映射后变得无法区分”编码进核。事实上
\[
\varphi(g_1)=\varphi(g_2)
\]
当且仅当
\[
\varphi(g_2^{-1}g_1)=e_H,
\]
也就是
\[
g_1\ker\varphi=g_2\ker\varphi.
\]
因此真正被 \(\varphi\) 区分的不是单个 \(g\)，而是核的陪集。

这直接得到第一同构定理。定义
\[
\widetilde\varphi:G/\ker\varphi
\longrightarrow \operatorname{Im}\varphi,
\]
\[
\widetilde\varphi(g\ker\varphi)=\varphi(g).
\]
上面的等价关系保证它不依赖陪集代表元；它还是同态，并且既单射又满射。因此
\begin{equation}
\boxed{
G/\ker\varphi\cong\operatorname{Im}\varphi
}.
\label{eq:first-isomorphism-theorem-restated}
\end{equation}

这个定理把商群从“形式定义”变成了自然结构：每当一个同态把某些元素压到同一个像上，商掉核就得到真正保留下来的群结构。

一个直接例子是
\[
\varphi:\mathbb Z\to U(1),
\qquad
n\mapsto \ee^{2\pi\ii n/m}.
\]
其核为 \(m\mathbb Z\)，像是 \(m\) 次单位根群，因此
\[
\mathbb Z/m\mathbb Z\cong \mu_m.
\]
这也是“模 \(m\) 加法”与“复平面上离散旋转”之间的精确同构。

接下来考察群作用与轨道--稳定子定理。

群不仅可以作为抽象乘法表存在，更常见的是它作用在某个集合 \(X\) 上。左作用是映射
\[
G\times X\to X,
\qquad (g,x)\mapsto g\cdot x,
\]
满足
\[
e\cdot x=x,
\qquad
(g_1g_2)\cdot x=g_1\cdot(g_2\cdot x).
\]
对固定 \(x\in X\)，轨道和稳定子定义为
\[
\operatorname{Orb}(x)=\{g\cdot x:g\in G\},
\]
\[
G_x=\{g\in G:g\cdot x=x\}.
\]

映射
\[
G/G_x\to\operatorname{Orb}(x),
\qquad
gG_x\mapsto g\cdot x
\]
是良定义双射，因此有限群情形有
\begin{equation}
\boxed{
|G|=|\operatorname{Orb}(x)|\,|G_x|
}.
\label{eq:orbit-stabilizer-restated}
\end{equation}
这就是轨道--稳定子定理。

正 \(n\) 边形的二面体群 \(D_n\) 给出最直观的例子。令 \(D_n\) 作用在顶点集合上。任一顶点的轨道包含全部 \(n\) 个顶点，而固定一个顶点的对称只有恒等与经过该顶点的相应反射，共 \(2\) 个。因此
\[
|D_n|=n\times2=2n.
\]
这个计数不是靠枚举得到，而是群作用结构的直接结果。

接下来考察一般有限群的 Cauchy 定理。

前面针对 Abel 群给出了归纳证明。利用群作用，可以证明不要求交换律的完整 Cauchy 定理：若素数 \(p\mid|G|\)，则 \(G\) 中存在阶为 \(p\) 的元素。

考虑集合
\[
X=\left\{(g_1,\ldots,g_p)\in G^p:
 g_1g_2\cdots g_p=e\right\}.
\]
前 \(p-1\) 个分量任取后，最后一个由
\[
g_p=(g_1\cdots g_{p-1})^{-1}
\]
唯一确定，所以
\[
|X|=|G|^{p-1}.
\]
因为 \(p\mid|G|\)，有
\[
p\mid|X|.
\]

循环群 \(C_p\) 通过循环移位作用在 \(X\) 上：
\[
(g_1,g_2,\ldots,g_p)
\mapsto
(g_2,\ldots,g_p,g_1).
\]
需要先确认移位后仍属于 \(X\)。若
\[
g_1g_2\cdots g_p=e,
\]
则
\[
g_2\cdots g_pg_1
=g_1^{-1}(g_1\cdots g_p)g_1=e.
\]
因此作用良定义。

由于作用群阶为素数 \(p\)，每个轨道大小只能是 \(1\) 或 \(p\)。固定点必须满足
\[
g_1=g_2=\cdots=g_p=g,
\]
并且
\[
g^p=e.
\]
固定点至少有 \((e,\\ldots,e)\) 一个。设固定点总数为 \(N\)。因为其余轨道大小都是 \(p\)，有
\[
|X|\equiv N\pmod p.
\]
又 \(p\mid|X|\)，所以
\[
N\equiv0\pmod p.
\]
既然 \(N\ge1\)，必有 \(N\ge p\)，因此除单位元外还存在 \(g\neq e\) 满足
\[
g^p=e.
\]
由于 \(p\) 为素数，\(O(g)\mid p\) 且 \(g\neq e\)，故
\[
\boxed{O(g)=p}.
\]
这证明了完整的 Cauchy 定理。Abel 情形的归纳证明仍然有价值，因为它直接服务于后面的有限 Abel 群分解；群作用证明则说明“交换性”并不是 Cauchy 定理本身所必需的。


接下来考察有限域乘法群循环性的另一条直接证明。

前面已经利用有限 Abel 群的结构说明了 $U_p=(\mathbb Z/p\mathbb Z)^\times$ 是循环群。这里再给一条更贴近有限域本身的证明，因为它清楚揭示了``多项式根的个数''为什么会强迫一个本原根存在，而且不需要预先知道 $U_p$ 的完整循环因子分解。

先证明一个有限 Abel 群中的小事实。设 $G$ 为有限 Abel 群，把所有元素阶的最小公倍数记为
\[
m=\operatorname{lcm}\{\operatorname{ord}(g):g\in G\}.
\]
把 $m$ 分解为
\[
m=\prod_{j=1}^r \ell_j^{a_j},
\]
其中 $\ell_j$ 两两不同。由最小公倍数的定义，对每个 $\ell_j^{a_j}$ 都能找到一个元素 $g_j\in G$，使得 $\operatorname{ord}(g_j)$ 含有这个最高的 $\ell_j$ 次幂。若
\[
\operatorname{ord}(g_j)=\ell_j^{a_j}q_j,
\qquad \gcd(q_j,\ell_j)=1,
\]
令
\[
h_j=g_j^{q_j}.
\]
则
\[
\operatorname{ord}(h_j)=\ell_j^{a_j}.
\]
因为 $G$ Abel，所有 $h_j$ 可交换；又因为这些阶两两互素，所以
\[
h=h_1h_2\cdots h_r
\]
的阶为各阶的乘积：
\[
\operatorname{ord}(h)=
\prod_{j=1}^r\ell_j^{a_j}=m.
\]
最后一句也可以直接验证。若 $h^s=e$，则
\[
h_1^s\cdots h_r^s=e.
\]
把等式提升到 $m/\ell_j^{a_j}$ 次幂，除 $h_j$ 外其余因子都被消掉，从而 $h_j^s=e$；所以每个 $\ell_j^{a_j}\mid s$，即 $m\mid s$。因此确有元素的阶等于所有元素阶的最小公倍数。

现在取
\[
G=\mathbb F_p^\times.
\]
记 $N=|G|=p-1$，并仍用 $m$ 表示所有元素阶的最小公倍数。Lagrange 定理说明每个元素的阶都整除 $N$，所以
\[
m\mid N,
\qquad m\le N.
\]
另一方面，按 $m$ 的定义，对所有 $a\in\mathbb F_p^\times$ 都有
\[
a^m=1.
\]
因此 $\mathbb F_p$ 中全部 $N=p-1$ 个非零元素都是多项式
\[
X^m-1
\]
的根。但域上的非零 $m$ 次多项式至多只有 $m$ 个根，所以
\[
N\le m.
\]
结合 $m\le N$ 得
\[
m=N=p-1.
\]
而前面已经证明 $G$ 中存在阶恰好为 $m$ 的元素，于是存在 $g\in\mathbb F_p^\times$ 满足
\[
\operatorname{ord}(g)=p-1.
\]
因此
\begin{equation}
\boxed{
\mathbb F_p^\times=\langle g\rangle\cong C_{p-1}
}.
\label{eq:finite-field-multiplicative-cyclic}
\end{equation}
这样的 $g$ 就叫模 $p$ 的本原根。

这个证明与正十七边形的关系非常直接。对 $p=17$，乘法群阶为
\[
16=2^4.
\]
取本原根 $g=3$ 后，指数群
\[
\langle3\rangle
\supset
\langle3^2\rangle
\supset
\langle3^4\rangle
\supset
\langle3^8\rangle
\supset
\{1\}
\]
的指数每次都是 $2$。Gauss 周期中对 16 个非零指数进行的``八个一组、四个一组、两个一组''，正是这条子群链在单位根指数上的陪集分解。后面的二次方程不是偶然碰巧出现，而是每一个指数为 $2$ 的子群扩张在周期和上留下的代数痕迹。

